\documentclass[11pt,a4paper]{article}

\usepackage[utf8]{inputenc}
\usepackage[T1]{fontenc}
\usepackage[ngerman]{babel}

\usepackage{amsmath,amssymb,amsthm,mathtools}
\usepackage[a4paper,margin=2.3cm]{geometry}
\usepackage{lmodern}
\usepackage{microtype}
\usepackage{hyperref}
\usepackage{enumitem}
\usepackage{tikz}
\usetikzlibrary{calc}

\title{Oktonische Hülle der Radikalzerlegung\\
	und Sonderstellungen von Primzahlkonstellationen}
\author{Kollaborative Skizze}
\date{\today}

\newtheorem{satz}{Satz}
\newtheorem{lemma}[satz]{Lemma}
\newtheorem{proposition}[satz]{Proposition}
\newtheorem{korollar}[satz]{Korollar}
\theoremstyle{definition}
\newtheorem{definition}[satz]{Definition}
\newtheorem{beispiel}[satz]{Beispiel}
\theoremstyle{remark}
\newtheorem{bemerkung}[satz]{Bemerkung}

\newcommand{\N}{\mathbb{N}}
\newcommand{\F}{\mathbb{F}_2}
\newcommand{\Ocal}{\mathbb{O}}
\newcommand{\ggT}{\operatorname{ggT}}
\newcommand{\Typ}{\operatorname{Typ}}
\newcommand{\vp}{v_p}
\newcommand{\sig}{\sigma}

\begin{document}
	\maketitle
	
	\begin{abstract}
		Wir heben die reduzierte Sieben-Typen-Klassifikation der geschichteten Radikalzerlegung
		auf eine formale Achterstruktur
		\[
		\Sigma \cong \F_2^3
		\]
		an und interpretieren diese als Indexsystem einer oktonischen Basis.
		Darauf aufbauend untersuchen wir die Sonderstellungen klassischer Primzahlkonstellationen:
		Primzahlzwillinge, Primzahldrillinge, Primzahlvierlinge und Primzahlfünflinge.
		Es zeigt sich, dass diese Konstellationen in der oktonischen Signaturwelt auf sehr wenige
		Richtungen konzentriert sind und insbesondere der Primzahlvierling eine ausgezeichnete
		neutral-projektive Stellung einnimmt.
	\end{abstract}
	
	\section{Formale Achterfamilie}
	
	Für Primzahlen \(p\ge 5\) betrachten wir die vier Familien modulo \(12\):
	\[
	E:=\{p:\ p\equiv 1 \pmod{12}\},\qquad
	A:=\{p:\ p\equiv 5 \pmod{12}\},
	\]
	\[
	B:=\{p:\ p\equiv 7 \pmod{12}\},\qquad
	C:=\{p:\ p\equiv 11 \pmod{12}\}.
	\]
	
	\begin{definition}[Formale Paritätssignatur]
		Sei \(R\) quadratfrei mit \(\ggT(R,6)=1\). Dann definieren wir
		\[
		\sig(R):=(\varepsilon_A(R),\varepsilon_B(R),\varepsilon_C(R))\in\F_2^3,
		\]
		wobei
		\[
		\varepsilon_A(R):=\#\{p\mid R:\ p\in A\}\bmod 2,
		\]
		\[
		\varepsilon_B(R):=\#\{p\mid R:\ p\in B\}\bmod 2,
		\qquad
		\varepsilon_C(R):=\#\{p\mid R:\ p\in C\}\bmod 2.
		\]
	\end{definition}
	
	Wir identifizieren die acht Elemente von \(\F_2^3\) mit den acht formalen Signaturen
	\[
	\Sigma:=\{\varnothing,A,B,C,AB,AC,BC,ABC\}
	\]
	via
	\[
	000\leftrightarrow \varnothing,\quad
	100\leftrightarrow A,\quad
	010\leftrightarrow B,\quad
	001\leftrightarrow C,
	\]
	\[
	110\leftrightarrow AB,\quad
	101\leftrightarrow AC,\quad
	011\leftrightarrow BC,\quad
	111\leftrightarrow ABC.
	\]
	
	Die bisherige reduzierte Sieben-Typen-Welt entsteht durch
	\[
	\varnothing \sim ABC \sim E.
	\]
	
	\section{Oktonische Hülle}
	
	\begin{definition}[Oktonische Basiszuordnung]
		Sei
		\[
		\{1,e_1,e_2,e_3,e_4,e_5,e_6,e_7\}
		\]
		eine Basis der Oktonionenalgebra \(\Ocal\).
		Wir setzen
		\[
		\varnothing \leftrightarrow 1,
		\]
		\[
		A\leftrightarrow e_1,\qquad
		B\leftrightarrow e_2,\qquad
		C\leftrightarrow e_3,
		\]
		\[
		AB\leftrightarrow e_4,\qquad
		AC\leftrightarrow e_5,\qquad
		BC\leftrightarrow e_6,\qquad
		ABC\leftrightarrow e_7.
		\]
	\end{definition}
	
	\begin{bemerkung}
		Dies ist zunächst eine \emph{Indexierung} der acht Basiskomponenten.
		Die volle oktonische Multiplikation verlangt zusätzlich eine orientierte
		Fano-Struktur mit Vorzeichenregel. Für die vorliegende Klassifikation genügt
		die formale Indexgeometrie.
	\end{bemerkung}
	
	\section{Fano-Geometrie der sieben nichttrivialen Signaturen}
	
	Die sieben nichtnulligen Elemente von \(\F_2^3\) bilden die Punkte der Fano-Ebene.
	Die Geraden entsprechen genau den Tripeln
	\[
	x,y,z\in \F_2^3\setminus\{0\},\qquad x+y+z=0.
	\]
	
	In unseren Bezeichnungen sind das:
	\[
	\{A,B,AB\},\quad \{A,C,AC\},\quad \{B,C,BC\},
	\]
	\[
	\{AB,AC,BC\},\quad \{A,BC,ABC\},\quad \{B,AC,ABC\},\quad \{C,AB,ABC\}.
	\]
	
	\begin{center}
		\begin{tikzpicture}[scale=1.15]
			\coordinate (A)   at (90:3);
			\coordinate (B)   at (18:3);
			\coordinate (AB)  at (-54:3);
			\coordinate (C)   at (-126:3);
			\coordinate (AC)  at (162:3);
			\coordinate (BC)  at (0,0);
			\coordinate (ABC) at (0,1.1);
			
			% Outer circle-ish points
			\foreach \P/\Name in {A/A,B/B,AB/AB,C/C,AC/AC,BC/BC,ABC/ABC}{
				\fill (\P) circle (2pt);
			}
			
			% Some Fano lines
			\draw (A) -- (B) -- (AB) -- cycle;
			\draw (A) -- (C);
			\draw (AC) -- (BC);
			\draw (B) -- (C);
			\draw (A) -- (BC);
			\draw (B) -- (AC);
			\draw (C) -- (AB);
			\draw (ABC) circle (1.9);
			
			\node[above]      at (A)   {$A$};
			\node[right]      at (B)   {$B$};
			\node[below right]at (AB)  {$AB$};
			\node[below left] at (C)   {$C$};
			\node[left]       at (AC)  {$AC$};
			\node[below]      at (BC)  {$BC$};
			\node[above]      at (ABC) {$ABC$};
		\end{tikzpicture}
	\end{center}
	
	\begin{bemerkung}
		Die Zeichnung ist hier nur als Inzidenzbild gemeint. Für eine echte oktonische
		Multiplikation müsste zusätzlich eine Orientierung der Geraden festgelegt werden.
	\end{bemerkung}
	
	\section{Primzahlkonstellationen im oktonischen System}
	
	Wir untersuchen nun die formale Signatur des Radikals einer Konstellation.
	Dabei werden die Primzahlen \(2\) und \(3\) als \emph{außerhalb} des \(E/A/B/C\)-Systems
	behandelt; falls \(3\) in der Konstellation vorkommt, wird sie als singulärer
	Randanker abgespalten.
	
	\subsection{Primzahlzwillinge}
	
	Ein Primzahlzwilling ist ein Paar \((p,p+2)\) aus Primzahlen.
	
	\begin{proposition}[Zwillinge]
		Sei \((p,p+2)\) ein Primzahlzwilling mit \(p>3\). Dann gibt es genau zwei
		oktonische Signaturmöglichkeiten:
		\[
		\sig\bigl((p,p+2)\bigr)=AB \quad\text{oder}\quad \sig\bigl((p,p+2)\bigr)=C.
		\]
		Genauer gilt:
		\[
		p\equiv 5 \pmod{12}\ \Longrightarrow\ (p,p+2)\sim (A,B)\ \Longrightarrow\ \sig=AB,
		\]
		\[
		p\equiv 11 \pmod{12}\ \Longrightarrow\ (p,p+2)\sim (C,E)\ \Longrightarrow\ \sig=C.
		\]
	\end{proposition}
	
	\begin{proof}
		Für \(p>3\) muss ein Primzahlzwilling die Form \((6n-1,6n+1)\) haben.
		Modulo \(12\) gibt es nur zwei Möglichkeiten:
		\[
		(5,7)\quad\text{oder}\quad (11,1).
		\]
		Im ersten Fall erhält man die Familien \(A,B\), also Signatur \(AB\).
		Im zweiten Fall erhält man \(C,E\), wobei \(E\) paritätisch neutral ist;
		also bleibt die Signatur \(C\).
	\end{proof}
	
	\begin{bemerkung}
		Der Ausnahmezwilling \((3,5)\) ist ein Randfall: Nach Abspaltung von \(3\) bleibt
		nur der Typ \(A\).
	\end{bemerkung}
	
	\subsection{Primzahldrillinge}
	
	Unter Primzahldrillingen verstehen wir hier die dichten Muster
	\[
	(p,p+2,p+6)\qquad\text{oder}\qquad (p,p+4,p+6).
	\]
	
	\begin{proposition}[Drillinge]
		Jeder Primzahldrilling enthält notwendig die Primzahl \(3\). Nach Abspaltung von \(3\)
		bleiben genau zwei Signaturtypen übrig:
		\[
		(3,5,7)\ \Longrightarrow\ \sig=AB,
		\]
		\[
		(3,7,11)\ \Longrightarrow\ \sig=BC.
		\]
	\end{proposition}
	
	\begin{proof}
		In drei ungeraden Zahlen mit Abstandsmuster \(0,2,6\) oder \(0,4,6\) liegt stets eine
		Zahl in der Klasse \(0 \pmod 3\). Soll die ganze Konstellation aus Primzahlen bestehen,
		so muss diese Zahl gleich \(3\) sein. Das liefert genau die beiden bekannten Drillinge
		\[
		(3,5,7),\qquad (3,7,11).
		\]
		Nach Abspaltung von \(3\) bleiben die Familien
		\[
		(A,B)\quad\text{bzw.}\quad (B,C),
		\]
		also die Signaturen \(AB\) bzw. \(BC\).
	\end{proof}
	
	\begin{bemerkung}
		Der Primzahldrilling ist im oktonischen System kein „freies“ Objekt des
		\(E/A/B/C\)-Raumes, sondern ein \emph{3-verankerter Randzustand}.
	\end{bemerkung}
	
	\subsection{Primzahlvierlinge}
	
	Der klassische Primzahlvierling hat die Form
	\[
	(p,p+2,p+6,p+8).
	\]
	
	\begin{proposition}[Vierlinge]
		Sei \((p,p+2,p+6,p+8)\) ein Primzahlvierling mit \(p>3\).
		Dann ist seine formale oktonische Signatur immer
		\[
		\sig = ABC.
		\]
		Unter der reduzierten Projektion fällt der Vierling also immer auf den neutralen Typ
		\[
		E.
		\]
	\end{proposition}
	
	\begin{proof}
		Für \(p>3\) sind modulo \(12\) nur zwei Familienmuster möglich:
		\[
		(5,7,11,1)\quad\text{oder}\quad (11,1,5,7).
		\]
		Das entspricht
		\[
		(A,B,C,E)\quad\text{oder}\quad (C,E,A,B).
		\]
		In beiden Fällen treten \(A,B,C\) jeweils genau einmal auf, während \(E\) neutral ist.
		Also ist die formale Signatur stets
		\[
		111=ABC.
		\]
		Nach der Reduktion \(ABC\sim E\) landet der Vierling im Typ \(E\).
	\end{proof}
	
	\begin{bemerkung}
		Das ist die bemerkenswerteste Sonderstellung: Der Primzahlvierling liegt formal auf der
		\emph{verborgenen achten Richtung} \(ABC\cong e_7\), erscheint reduziert aber als neutral.
		Im oktonischen Bild ist er also ein \emph{verdeckter Nichtnullzustand}, der in der
		Siebentypenprojektion zu \(E\) kollabiert.
	\end{bemerkung}
	
	\subsection{Primzahlfünflinge}
	
	Der dichte Primzahlfünfling hat die Form
	\[
	(p,p+2,p+6,p+8,p+12).
	\]
	
	\begin{proposition}[Fünflinge]
		Sei \((p,p+2,p+6,p+8,p+12)\) ein Primzahlfünfling mit \(p>3\).
		Dann gibt es genau zwei formale oktonische Signaturen:
		\[
		p\equiv 5 \pmod{12}\ \Longrightarrow\ \sig = BC,
		\]
		\[
		p\equiv 11 \pmod{12}\ \Longrightarrow\ \sig = AB.
		\]
	\end{proposition}
	
	\begin{proof}
		Für \(p>3\) gibt es modulo \(12\) genau zwei Muster:
		\[
		(5,7,11,1,5)\quad\text{oder}\quad (11,1,5,7,11).
		\]
		Das entspricht
		\[
		(A,B,C,E,A)\quad\text{oder}\quad (C,E,A,B,C).
		\]
		Im ersten Fall ist die Parität
		\[
		A:2,\quad B:1,\quad C:1,
		\]
		also reduziert
		\[
		(0,1,1)=BC.
		\]
		Im zweiten Fall ist die Parität
		\[
		A:1,\quad B:1,\quad C:2,
		\]
		also reduziert
		\[
		(1,1,0)=AB.
		\]
	\end{proof}
	
	\begin{bemerkung}
		Im Gegensatz zum Vierling projiziert der Fünfling nicht neutral, sondern spaltet sich
		in zwei echte oktonische Richtungen auf.
	\end{bemerkung}
	
	\section{Zusammenfassung der Sonderstellungen}
	
	\begin{satz}[Sonderlagen klassischer Konstellationen]
		Im formalen oktonischen Signatursystem besitzen die klassischen Primzahlkonstellationen
		folgende ausgezeichnete Lagen:
		\begin{enumerate}[label=\textup{(\roman*)}]
			\item Primzahlzwillinge mit Abstand \(2\) liegen nur in den Richtungen
			\[
			C \quad\text{oder}\quad AB.
			\]
			\item Primzahldrillinge sind notwendig \(3\)-verankert und liefern nach Abspaltung von \(3\)
			nur die Richtungen
			\[
			AB \quad\text{oder}\quad BC.
			\]
			\item Primzahlvierlinge liegen formal immer in der Richtung
			\[
			ABC,
			\]
			projizieren reduziert aber auf
			\[
			E.
			\]
			\item Primzahlfünflinge liegen nur in den Richtungen
			\[
			AB \quad\text{oder}\quad BC.
			\]
		\end{enumerate}
	\end{satz}
	
	\begin{proof}
		Die vier Aussagen sind die vorstehenden Propositionsresultate.
	\end{proof}
	
	\section{Interpretation}
	
	\begin{bemerkung}
		Das oktonische Bild legt folgende Lesart nahe:
		\begin{itemize}[itemsep=0.2em]
			\item Zwillinge sind \emph{zweirichtige Randmodi}.
			\item Drillinge sind \emph{durch \(3\) erzwungene singuläre Konfigurationen}.
			\item Vierlinge sind \emph{oktonisch verdeckte Neutralzustände}: formal \(ABC\), reduziert \(E\).
			\item Fünflinge verlassen diese Neutralität wieder und spalten sich in zwei echte
			nichtneutrale Richtungen auf.
		\end{itemize}
	\end{bemerkung}
	
	\begin{bemerkung}
		Gerade die Sonderrolle des Vierlings ist mathematisch interessant: Er ist die erste dichte
		Konstellation, deren volle Familienparität exakt alle drei aktiven Richtungen \(A,B,C\)
		umfasst und dadurch auf die \(ABC\)-Achse fällt.
	\end{bemerkung}
	
	\section{Schluss}
	
	Die oktonische Hülle ist keine neue Primzahlsatz-Theorie, aber sie organisiert die
	Familienparitäten der Radikalschichten in einer erstaunlich straffen Geometrie.
	Die klassischen Primzahlkonstellationen besetzen darin nur sehr wenige ausgezeichnete
	Richtungen, und insbesondere der Primzahlvierling nimmt eine singuläre Zwischenstellung
	zwischen verborgener Achterstruktur und reduzierter Siebentypenprojektion ein.
	
\end{document}